\chapter{Soft Prompting as a Residual-Stream Intervention}
\label{sec:k_e}

\paragraph{Continuous Prefix Parameterisation.}
We use soft prompting as a control intervention experiment, to test whether a small 
input-side perturbation can improve first-token arithmetic or judgement behaviour,
 and whether that perturbation is interpretable relative to the task. 

Let $x=(t_1,\ldots,t_T)$ be a tokenised prompt and let $E\in\mathbb{R}^{V\times d}$ 
be the embedding matrix. A soft prompt learns $P\in\mathbb{R}^{k\times d}$ and 
prepends it to the real token embeddings,
\begin{equation}
\bigl(P_1,\ldots,P_k,e_{t_1},\ldots,e_{t_T}\bigr)
\in \mathbb{R}^{(k+T)\times d}.
\label{eq:soft_prompt_seq}
\end{equation}
We prepend $k$ pad-token positions to the input and replace them with the rows of $P$
at the embedding layer before the forward pass, so the prefix enters the residual
stream while all base-model parameters remain frozen and only $P$ is updated during training.

\paragraph{Training Objective.}

The objective is cross-entropy on the first answer token at the final prompt position,
\begin{equation}
\mathcal{L}
=
-\frac{1}{B}\sum_{b=1}^{B}
\log
\frac{\exp f_{y_b}(x_b)}
     {\sum_v \exp f_v(x_b)}.
\label{eq:soft_prompt_loss}
\end{equation}
We train with AdamW ($\mathrm{lr}=10^{-3}$, weight decay $10^{-4}$), 
cosine annealing, and unit-norm gradient clipping on a $90/10$ train--validation
 split, retaining the lowest validation-loss checkpoint.

\paragraph{Results.}

We evaluate one prefix length, $k=10$, on three tasks. The residue-class task uses prompts of the form \promptinline{calc: \{a\}\%7=}, where the model must predict the single-digit remainder modulo~7. The balanced-parentheses and causal-direction tasks use binary yes/no prompts of the form \promptinline{Is \{seq\} a valid bracket string? Answer yes or no:} and \promptinline{Does \{A\} lead to \{B\}? Answer yes or no:}, and test whether the prefix can correct a response-format failure.

\begin{table}[ht!]
\centering
\small
\begin{tabular}{@{}lrrrr@{}}
\toprule
Task & Base acc.\ (\%) & Prefix acc.\ (\%) & $\Delta$ acc.\ (\%) & $\Delta\,\bar{p}(\mathrm{correct})$ \\
\midrule
Residue class        & 15.0 & 16.0 & $+1.0$  & $-0.011$ \\
Balanced parentheses &  0.0 & 28.6 & $+28.6$ & $+0.078$ \\
Causal direction     &  0.0 & 53.8 & $+53.8$ & $+0.236$ \\
\bottomrule
\end{tabular}
\caption{Soft-prompt validation results for Qwen3-4B with frozen base weights and prefix length $k=10$.}
\label{tab:sp_summary}
\end{table}

Table~\ref{tab:sp_summary} shows that for residue-class completion 
the prefix does not improve performance, 
as the mean probability of the correct 
token decreases despite a negligible accuracy change, 
suggesting the model follows a leading-digit heuristic 
that the prefix weakens without substituting a reliable 
modular computation. For balanced parentheses and causal
 direction, the baseline consistently predicts the token
  \texttt{1}, so the primary failure is response format 
  rather than task knowledge, and the prefix redirects 
  probability mass toward yes/no tokens. We conclude that
   a learned prefix can correct an output-format 
   failure when the model already encodes the relevant 
   computation, but cannot introduce a computation that
    is absent, placing soft prompting here in a diagnostic
     rather than a knowledge-editing role.
